\newpage
\section{两个重要的常数}
\label{sec:analysis-constants}

\subsection{自然对数的底\texorpdfstring{\(e\)}{e}}
\(e\)是通过数列\(\left\{ e_{n} \right\}\)的极限引进的一个常数，\(e \approx 2.71828\)，其中\(e_{n} = {\left( 1 + \frac{1}{n} \right)}^{n}\).

\begin{example}
已知正数\(a\)，将其分成若干部分，若要使它们的乘积达到最大，应该怎样分？
\end{example}

\subsection{关于数\texorpdfstring{\(e\)}{e}的几个基本命题}
设\(e_{n} = {\left( 1 + \frac{1}{n} \right)}^{n} , n \in \mathbb{N}^{*}\)，证明：\(\left\{ e_{n} \right\}\)严格单调递增且收敛.

（1）证明：\(e = \sum _{n = 0}^{\infty} \frac{1}{n !}\).

（2）记\(\varepsilon_{n} = e - \sum _{k = 0}^{n} \frac{1}{k !}\)，证明：\(\lim_{n \to \infty} \varepsilon_{n} \left( n + 1 \right) ! = 1\).

（3）对于上述\(\varepsilon_{n}\)有\(\frac{1}{\left( n + 1 \right) !} < \varepsilon_{n} < \frac{1}{n ! n}\).

（4）\(e\)是无理数.

（5）\(\lim_{n \to \infty} \frac{n}{\sqrt[n]{n !}} = e\).

\subsection{Euler常数\texorpdfstring{\(\gamma\)}{gamma}}
\label{subsec:analysis-euler-constant}
证明：数列\(\left\{ \gamma_{n} \right\}\)收敛，其中\(\gamma_{n} = \sum _{i = 1}^{n} \frac{1}{i} - \ln n , n \in \mathbb{N}^{*}\)。其极限称为{\heiti{}Euler常数}，记为\(\gamma\)。

计算极限：
\begin{equation}
    \lim_{n \to \infty} \left( \frac{1}{n + 1} + \frac{1}{n + 2} + \cdots + \frac{1}{2 n} \right)
\end{equation}

{\kaiti{}思考\quad Euler常数\(\gamma\)是否是无理数？}

\begin{example}
    计算以下极限：

    （1）\(\lim_{n \to \infty} {\left( 1 - \frac{1}{n} \right)}^{n}\),
    \(\lim_{n \to \infty} {\left( 1 + \frac{1}{2 n} \right)}^{n}\)

    （2）\(\lim_{n \to \infty} {\left( 1 + \frac{1}{n} \right)}^{n^{2}}\)

    （3）\(\lim_{n \to \infty} {\left( 1 + \frac{1}{n} + \frac{1}{n^{2}}
        \right)}^{n}\)

    （4）\(\lim_{n \to \infty} \prod _{i = 1}^{n} \left( 1 +
    \frac{i}{n^{2}} \right)\)

    （5）\(\lim_{n \to \infty} \frac{n ! 2^{n}}{n^{n}}\)

    （6）\(\lim_{n \to \infty} {\left( 1 +
        \frac{1}{p_{n}}\right)}^{p_{n}}\)，其中\(\left\{ p_{n} \right\}\)是正数列
\end{example}

